\chapter*{Conclusion}
	\graphicspath{{Conclusion/}}
	\addcontentsline{toc}{chapter}{Conclusion}
	\label{chap:conclusion}

Cette thèse avait pour objectif d'étudier différentes méthodes de reconstruction tomographique afin de les appliquer dans le cas de l'imagerie dentaire à faible dose. Pour cela, nous avons adapté des méthodes itératives utilisant une régularisation avec variation totale aux spécificités de la géométrie conique majoritairement utilisée en dentaire, et nous avons évalué les performances de réseaux de neurones en post-traitement de reconstructions analytiques. 

Nous avons commencé par introduire les principes physiques et mathématiques de la reconstruction tomographique, tout en insistant sur les caractéristiques de l'imagerie dentaire CBCT. De premières méthodes de reconstruction ont été présentées~: les reconstructions analytiques, rapides et efficaces en conditions idéales, mais sensibles au bruit et aux projections en faible nombre\, ; et certaines reconstructions itératives permettant de prendre en compte la distribution du bruit dans les projections. 

Dans un second temps, nous avons étudié les algorithmes de reconstruction tomographique itératifs incluant une régularisation TV, adaptés à une distribution gaussienne ou Poisson pour le bruit. Nous les avons appliqués au fantôme de Shepp-Logan 2D en étudiant leur performance sur un bruit mixte gaussien-Poisson. Ces résultats préliminaires ont permis de mettre en avant l'importance du choix du paramètre de régularisation dans le résultat de la reconstruction. 

Ces algorithmes ont ensuite été adaptés et appliqués progressivement à la géométrie conique avec les fantômes de Shepp-Logan et de mâchoire, puis à des données expérimentales acquises sur un fantôme anthropomorphique. Ces expériences ont mis en valeur les avantages et les inconvénients de ces méthodes. Elles permettent d'un côté de réduire efficacement les artefacts et le bruit présent dans les reconstructions analytiques faibles-doses, avec un compromis important entre conservation des détails et réduction du bruit, compromis géré par le paramètre de régularisation. D'un autre côté, ces méthodes nécessitent d'importantes ressources de calcul numérique, en particulier en CBCT et dans le cas de projections tronquées, et les reconstructions sont relativement lentes pour une utilisation clinique. De plus, la régularisation TV utilisée dans ces méthodes n'est pas forcément la plus adaptée à l'imagerie dentaire~: elle favorise les zones constantes par morceaux et est assez efficace quand il s'agit de reconstruire de grosses structures telles que les dents, mais a tendance à effacer les détails particulièrement fins (canaux de dents, structure d'os). Cette étude a également mis en avant la difficulté de comparer les différentes méthodes itératives régularisées entre elles. La fonction coût différente entre SIRT-TV et les deux autres, ou encore la différence d'architecture entre l'algorithme KL-TV et les autres rendent délicat l'obtention d'une comparaison juste. 

Nous pouvons lister plusieurs manières d'améliorer le travail effectué sur les méthodes itératives. Tout d'abord, il serait possible d'accélérer les calculs en utilisant OSEM \citepp{Hudson1994}, ou l'algorithme de Chambolle-Pock stochastique \citepp{Chambolle2018} qui n'utilisent qu'une partie des projections à chaque itération. L'ensemble des projections choisies change à chaque itération de sorte à ce que la totalité des projections soit prise en compte régulièrement. Ces méthodes pourraient réduire les temps de calcul et la mémoire nécessaire. Du point de vue de la régularisation, la variation totale étudiée ici, n'est pas adaptée pour reconstruire les structures très fines. Les reconstructions itératives régularisées par TV souffrent ainsi de la perte de petits détails tels que la structure des os, ou certains canaux. Il serait donc intéressant d'utiliser d'autres régularisations. \cite{Ye2018} ont proposé un algorithme itératif dans lequel la régularisation est effectuée par un réseau de neurones.  

\paragraph{} Un deuxième axe d'étude d'amélioration de la reconstruction tomographique a été étudié dans cette thèse~: l'application des méthodes d'apprentissage profond, particulièrement efficaces en traitement d'images depuis quelques années. Nous avons synthétisé quelques travaux présents dans l'état de l'art, en expliquant qu'elles devaient être adaptées en fonction du type de données à traiter. Ainsi, disposant d'une base de données de projections acquises en dose normale, nous avons décidé d'utiliser des méthodes d'apprentissage supervisé, en simulant des acquisitions faibles doses à partir des données à notre disposition. Nous avons pris une architecture relativement simple mais qui a déjà prouvé son efficacité dans de nombreux domaines d'application~: les U-Net. Plusieurs U-Net ont été entraînés à réduire les artefacts présents dans les reconstructions analytiques effectuées à partir des données faibles-doses. Nos expériences ont montré de très bons résultats pour le U-Net multi-plans, simple et relativement rapide à entraîner, aboutissant à des reconstructions précises et de qualité, permettant d'effectuer plus facilement de meilleurs diagnostics médicaux.

Les bonnes performances des réseaux de neurones doivent tout de même être mises en perspective. La phase d'entraînement est cruciale dans la capacité du réseau à généraliser ce qu'il a appris sur de nouvelles données. Notre base de données, constituée de 32 patients, a déjà montré quelques exemples de cas où la sous-représentation des données a pu poser des problèmes. Pour un usage clinique de ces algorithmes, il est donc nécessaire de les ré-entraîner avec une plus grande base de données. 

Les expériences menées dans cette thèse donnent un aperçu des bonnes capacités des méthodes d'apprentissage profond à améliorer l'aspect de reconstructions analytiques. Ces méthodes sont néanmoins justement basées sur une reconstruction analytique. Dans nos simulations de faible dose, nous avons choisi un nombre de projections suffisamment élevé pour que les reconstructions analytiques restent correctes, malgré la corruption par le bruit et les artefacts. Nous pouvons facilement imaginer que dans des cas plus extrêmes, la reconstruction analytique ne permette pas de récupérer assez d'informations pour qu'un réseau puisse restaurer correctement les images par la suite. Une façon d'éviter ce problème serait d'inclure le processus de reconstruction au réseau de neurones, en utilisant les réseaux de unrolling. La nécessité de très grandes capacités numériques, plus particulièrement en cone beam, fait qu'en pratique, ces algorithmes sont encore difficiles à mettre en \oe uvre. De récents articles laissent cependant espérer une possible utilisation des méthodes, en considérant par exemple une approche de type OSEM dans lequel les projections ne seraient pas toutes utilisées simultanément \citepp{Tang2021} ; ou encore en utilisant une version sous-échantillonnée des volumes et des opérateurs de projections/rétro-projection \citepp{Tang2022}.

Nous rappelons également qu'une large variété de réseau, est aujourd'hui disponible et que cette thèse n'a étudié que des U-Net. L'entraînement d'autres réseaux, tels que les GAN ou les \textit{Transformers}, pourrait étayer cette étude sur les capacités des réseaux de neurones à améliorer la qualité d'images CT faibles-doses. Enfin, nous pourrions effectuer d'autres tâches simultanément au débruitage, et par exemple prendre en compte le flou dans les images et le corriger avec de la déconvolution.

\newpage
\bibliographystyle{apalike-fr}
\bibliography{bib}









